\documentclass[a4paper, 11pt]{article}


\usepackage[slovak]{babel}
\usepackage[utf8]{inputenc}
\usepackage[left=2cm, top=3cm, text={17cm, 24cm}]{geometry}
\usepackage{times}
\usepackage{verbatim}
\usepackage{enumitem}
\usepackage{graphicx} 
\usepackage[unicode]{hyperref}
\usepackage{listings}
\usepackage{color}


\begin{document}

	\begin{titlepage}
		\begin{center}
			\Huge
			\textsc{Vysoké učení technické v~Brně} \\
			\huge
			\textsc{Fakulta informačních technologií} \\
			\vspace{\stretch{0.382}}
			\begin{figure}[h]
                \centering
                \scalebox{0.6}{\includegraphics{logo.eps}}
            \end{figure}
			\LARGE
			Projektová dokumentácia k~predmetu IMS \\
			\Huge
			Zadanie č. 3 - Epidemiologické modely - mikro-úroveň
			\vspace{\stretch{0.618}}
		\end{center}

		{\Large
			\today
			\hfill
			\begin{tabular}{r}
			Dávid Oravec (xorave05) \\
			Dominik Boboš (xbobos00)
			\end{tabular}
		}
	\end{titlepage}
	\newpage
	
	\tableofcontents
	\newpage
	
	\section{Úvod}
	\par Pre vytvorenie epidemiologického modelu, sme si vybrali ochorenie COVID-19. Je to infekčné ochorenie vyvolané koronavírusom SARS-CoV-2, ktorý patrí do rodu Betacoronavirus. Prvý oficiálne potrvdený prípad sa datuje na 8. december 2019, aj keď niektoré nepotvrdené zdroje uvádzajú 17. november 2019 \cite{guardian}.
	\newline
	
	\par Ochorenie postihlo celý svet a tak po roku od prvého prípadu sa stále nachádza v pandemickej kríze. COVID-19 zasiahol do života ľudí a v boji proti tejto pandémii musia vlády všetkých krajín zavádzať opatrenia, ktoré pomôžu zabrániť šíreniu. Väčšina týchto opatrení má devastačné následky v oblasti ekonomiky, kultúry, športu, ale aj psychického zdravia ľudí.
	\newline
	
	\par V súvislosti s touto problematikou sme sa rozhodli zamerať na skúmanie šírenia koronavírusu v oblasti kultúry, konrétne kinematografie. Našim cieľom je zistiť, akým spôsobom sa vírus môže šíriť medzi ľuďmi, ktorí sú uzavretí v realtívne malom priestore (kinosále) a ako znížiť riziko nákazy bez potreby zavádzania drastických opatrení.
	
    	\subsection{Zdroje faktov}
    	    \subsubsection{Autori}
    	    \begin{itemize}
    	        \item Dominik Boboš (\href{mailto:xbobos00@stud.fit.vutbr.cz}{xbobos00@stud.fit.vutbr.cz})
    	        \item Dávid Oravec (\href{mailto:xorave05@stud.fit.vutbr.cz}{xorave05@stud.fit.vutbr.cz})
    	    \end{itemize}
    	    
    	    \subsubsection{Literatúra}
    	    Ako primárne zdroje pre našu štúdiu a zároveň \uv{odrazový mostík} sme použili \cite{maxbrenner2020}, \cite{charlesconnor2020} a \cite{joslelieveldfrankhelleiset.al.2020}. Ostatné literárne zdroje sú uvedené v sekcii \hyperref[sec:bib]{Bibliografia}.
	    
	\section{Rozbor témy}
	\par COVID-19 je relatívne nový vírus a preto o ňom vieme stále málo. Jeho špecifickosť spočíva v tom, že sa neprejavuje u každého človeka rovnakým spôsobom. Inkubačná doba vírusu je od 2 do 14 dní, pričom infekciu je možné prenášať aj počas inkubačnej doby (v období bez príznakov ochorenia). Spôsob prenášania vírusu je prostredníctvom kvapôčok, ktoré sú vylučované pri kašli, kýchaní či rozprávaní. V ohrození sú najmä ľudia, ktorí sú v kontakte s nakazeným dlhšie ako 15 minút a ich odstup je menší než 2 metre \cite{wikipedia}. 
	
	\par V tejto štúdii sa zameriavame na modelovanie šírenia nákazy v kinosále počas filmového predstavenia. Na základe štúdie \cite{joslelieveldfrankhelleiset.al.2020} sme zostavili potrebné rovnice, ktoré pomôžu monitorovať priebeh nákazy medzi návštevníkmi kina. Pre náš model je využitý celulárny automat.
	
	Možno tu ešte niečo dopísať I am Lost
	
	    \subsection{Celulárny automat}
        Na experimenty je využitý celulárny automat (CA), pomocou ktorého sledujeme šírenie nákazy medzi jednotlivými bunkami (návštevníkmi kina). CA je definovaný niekoľkými prvkami \cite{petrperingermartinhruby}.	
        \begin{itemize}
            \item \textbf{Bunka} - základný element, môže byť v jednom z konečného počtu stavov
            \item \textbf{Pole buniek} - n-rozmerné, rovnomerné rozdelenie priestoru
            \item \textbf{Okolie} - počet a pozícia okolitých buniek s ktorými sa pracuje
            \item \textbf{Pravidlá} - stav bunky a jej okolia definujúci nový stav bunky
        \end{itemize}
        Jednotlivé generácie buniek sú vymodelované a modifikované práve vďaka funkcionalite CA. 
    
    \section{Koncepcia modelu}
    V našom modeli skúmame šírenie vírusu za rôznych podmienok. Podmienky (opatrenia) záležia na uživateľskom vstupe. Náš model podporuje opatrenia v kinosálach: 
        \begin{itemize}
            \item návštevníci majú povinne nasadené rúška
            \item v kinosále je striedavé sedenie (1 voľná sedačka medzi návštevníkmi sediacimi v jednej rade)
            \item v kine je zapnutá ventilácia vzduchu
        \end{itemize}
    Na základe týchto opatrení skúmame šírenie vírusu medzi návštevníkmi počas filmu. Predpokladáme vždy, že je sála vypredaná a teda všetky miesta, na ktorých je práve umožnené sedieť, sú obsadené.
    
        \subsubsection{Stavy bunky}
        Bunky môžu byť v jednom z následujúcich stavov \{-1,0,1,2\}.
        \begin{itemize}
            \item \textbf{-1} je stav, ktorý reprezentuje zakázané sedenie (pri opatrení na obmedzenie kapacity kinosály)
            \item \textbf{0} je stav, ktorým disponuje zdravý návštevník kina
            \item \textbf{1} je stav, ktorým disponuje nakazený návštevník kina
            \item \textbf{2} je stav, kedy je na zdravého návštevníka prenesený vírus počas pobytu v kinosále
            \begin{itemize}
                \item v tomto prípade však nakazený návštevník nie je hneď infekčný a neroznáša vírus
            \end{itemize}
        \end{itemize}
        
        \subsubsection{Pole buniek}
        Pole buniek reprezentuje kinosálu. Kapacita kinosály je 200 osôb, pri čom sedadlá sú rozdelené do 10 radov po 20 miest. Výnimkou, kedy má sála obmedzenú kapacitu, je pri zavedení opatrenia na striedavé sedenie. Vtedy je kapacita sály 100 miest. Tak ako už bolo spomenuté, predpokladáme zaplnenie sály vždy na 100\%.
        
        \subsubsection{Okolie}
        Model ráta s okolím bunky typu \texttt{Moore}, čiže 3x3. To znamená, že vírus sa môže šíriť od všetkých susediacich buniek \cite{petrperingermartinhruby}.
    
    
	\newpage
    \bibliographystyle{czechiso}
    \renewcommand{\refname}{Bibliografia}
    \label{sec:bib}
    \bibliography{ims}

\end{document}